\chapter{Ölçüm Düzeni ve Sonuçlar}
\label{ch:sonuclar}

\section{Ölçüm Planı}
\label{sec:olcum-plani}

Gelişme raporundaki 20--30 kişilik ölçüm hedefi, sergi takvimi içinde
gerçekleştirilemediği için pratik düzeyde \textbf{küçültülmüş bir test
seti} kullanılmıştır. Test seti üç bileşenden oluşur:

\begin{itemize}[itemsep=2pt]
    \item \textbf{Gallery (kayıtlı):} Proje ekibi (Göktürk Göçen, Ekin
          Ağaoğlu) ve gönüllü olarak katılan toplam 5--8 kişi. Her bir
          birey için kayıt fotoğrafı düzgün aydınlatma ve frontal poz
          ile alınmıştır.
    \item \textbf{Genuine probe seti:} Aynı bireylerin farklı zamanlarda,
          farklı poz/aydınlatmalarda alınmış 3--5 fotoğrafı.
    \item \textbf{Impostor probe seti:} Gallery'de olmayan bireylere ait
          (kamuya açık veri setlerinden seçilmiş) fotoğraflar.
\end{itemize}

Her probe \texttt{far\_frr.py} aracı ile 10 karelik bir burst olarak
sunucuya gönderilmiş; sunucu kalite filtresi, centroid hesabı ve kosinüs
benzerliği üretmiştir. Test üç farklı koşul grubunda tekrarlanmıştır:

\begin{enumerate}[label=(\arabic*),itemsep=2pt]
    \item \textbf{Aydınlatma}: 100~lx (alacakaranlık), 300~lx (ofis), 600~lx (gün ışığı).
    \item \textbf{Mesafe}: 30--60~cm (yakın), 60--90~cm (orta), 90--120~cm (uzak).
    \item \textbf{Pasif canlılık}: gerçek yüz vs.\ ekran tekrar (replay) testleri.
\end{enumerate}

\paragraph{Veri Yönetimi.} Tüm test görüntüleri sergi sonrası
silinmek üzere yalnızca yerel diskte saklanmıştır; gönüllü katılımcılardan
açık rıza alınmıştır. Sunucu tarafında ham görüntü diske yazılmamış,
yalnızca \texttt{far\_frr.py}'nin CSV çıktıları (yüz görüntüsü yok,
similarity skorları var) raporlamaya alınmıştır.

\section{Sonuç Tabloları (Yer Tutucu)}
\label{sec:sonuc-tablolari}

Ölçüm dönemi tamamlandıkça aşağıdaki tablolar gerçek değerlerle
doldurulacaktır. Mevcut sürümde değerler ``--'' ile bırakılmıştır.

\begin{table}[h!]
\centering
\caption{Aydınlatma koşullarına göre FAR/FRR (orta mesafe, $\tau$ = 0.40).}
\label{tab:results-light}
\small
\begin{tabular}{l c c c c}
\toprule
\textbf{Koşul} & \textbf{Genuine probe} & \textbf{Impostor probe} & \textbf{FAR} & \textbf{FRR}\\
\midrule
100~lx (alacakaranlık) & -- & -- & -- & --\\
300~lx (ofis)          & -- & -- & -- & --\\
600~lx (gün ışığı)     & -- & -- & -- & --\\
\bottomrule
\end{tabular}
\end{table}

\begin{table}[h!]
\centering
\caption{Mesafeye göre FAR/FRR (300~lx, $\tau$ = 0.40).}
\label{tab:results-distance}
\small
\begin{tabular}{l c c c c}
\toprule
\textbf{Mesafe}  & \textbf{Genuine probe} & \textbf{Impostor probe} & \textbf{FAR} & \textbf{FRR}\\
\midrule
30--60~cm   & -- & -- & -- & --\\
60--90~cm   & -- & -- & -- & --\\
90--120~cm  & -- & -- & -- & --\\
\bottomrule
\end{tabular}
\end{table}

\begin{table}[h!]
\centering
\caption{Pasif canlılık skoru — gerçek yüz vs.\ fotoğraf replay (300~lx, 60--90~cm).}
\label{tab:results-liveness}
\small
\begin{tabular}{l c c}
\toprule
\textbf{Koşul} & \textbf{Ortalama $\sigma_{emb}$} & \textbf{Standart sapma}\\
\midrule
Canlı yüz (gallery üyeleri) & -- & --\\
Ekran tekrarı (telefonda gösterilen foto) & -- & --\\
\bottomrule
\end{tabular}
\end{table}

\section{Operasyonel Performans}
\label{sec:perf}

Sistemin uçtan uca gecikmesi, \texttt{far\_frr.py} aracının her burst için
ölçtüğü değerlerden hesaplanmıştır. \Cref{tab:results-latency} kategori
bazında ortalama, p50, p95 ve maksimum gecikme değerlerini gösterir.

\begin{table}[h!]
\centering
\caption{Uçtan uca gecikme dağılımı (n adet burst üzerinden, saniye).}
\label{tab:results-latency}
\small
\begin{tabular}{l c c c c c}
\toprule
\textbf{Kategori} & \textbf{n} & \textbf{Ortalama} & \textbf{p50} & \textbf{p95} & \textbf{Maks}\\
\midrule
Tüm burst'ler                          & -- & -- & -- & -- & --\\
Genuine eşleşme                         & -- & -- & -- & -- & --\\
Impostor reddi                          & -- & -- & -- & -- & --\\
Kalite yetersiz (insufficient frames)   & -- & -- & -- & -- & --\\
\bottomrule
\end{tabular}
\end{table}

Gecikme bileşenlerinin yaklaşık dağılımı (varsayılan koşulda, hedef):

\begin{itemize}[itemsep=1pt]
    \item ESP32-CAM görüntü yakalama (10 kare $\times$ 200~ms): $\sim$2.0~s
    \item HTTP POST (Wi-Fi $\geq$~10\,Mbps, 10 kare): $\sim$1.5~s
    \item Sunucu InsightFace çıkarımı (10 kare, CPU): $\sim$2.0~s
    \item Centroid + DB araması: $<$~50~ms
    \item HTTP cevap dönüşü + UART iletim: $<$~200~ms
    \item \textbf{Beklenen toplam:} 5--7~s
\end{itemize}

\section{Eşik Optimizasyonu}

\Cref{ch:gerceklestirme}'de tanımlanan \texttt{far\_frr.py} aracı,
$\tau \in [0, 1]$ üzerinde tam süpürme yapar. EER ve $\mathrm{FAR} \leq 0.01$
kısıtı altındaki operating point, ölçüm tamamlandıktan sonra
aşağıdaki gibi raporlanacaktır:

\begin{itemize}[itemsep=1pt]
    \item \textbf{EER eşiği:} $\tau^\star_{EER}$ = -- (FAR = FRR = --)
    \item \textbf{Operatif eşik} ($\mathrm{FAR} \leq 0.01$): $\tau^\star_{op}$ = -- (FRR = --)
\end{itemize}

Sunucu \texttt{MATCH\_THRESHOLD} ortam değişkeni, ölçüm sonrası optimize
edilen değere göre güncellenip Docker konteyneri yeniden başlatılır.

\section{Donanım Entegrasyon Doğrulaması}

\paragraph{ESP32-CAM $\rightarrow$ EC2.} Uçtan uca burst akışı tezgahta
doğrulanmıştır: ESP32-CAM Wi-Fi STA modunda hotspot'a bağlanmakta, 10
kareyi sunucuya POST etmekte ve JSON yanıtı parse etmektedir. Server
tarafındaki yapılandırılmış JSON günlüğü, gelen burst sayısını,
kalite-geçen kare sayısını ve gecikmeyi izlemekte kullanılmıştır.

\paragraph{STM32 saat ve UART köprüsü.} HSI$\times$PLL ile 216~MHz HCLK
üretilmiş, USART3 VCP'den \texttt{printf} debug çıktısı ve USART6
üzerinden ESP32-CAM'e \texttt{CAPTURE/SCANNING/MATCH/\ldots} satırları
doğrulanmıştır. State machine modülü host-side bir test harness'ında
12 senaryo için ayrıca deterministik olarak doğrulanmıştır
(\Cref{tab:host-tests}).

\paragraph{ESP32-CAM BLE çevresel.} ESP32-CAM cihazı, BLE peripheral
olarak \texttt{TaxiGuard} adıyla yayın yapmakta; \texttt{FFE0} servis
ve \texttt{FFE1} karakteristik üzerinden \emph{notify} akışı
çalışmaktadır. MTU 247 yapılandırması, \texttt{MATCH:<name>;<sim>}
gibi 20 baytı aşan satırların tek pakette iletilmesini garanti eder.

\paragraph{iPhone (TaksiGuvenlik).} SwiftUI uygulaması Login/Register
akışından geçirilmiş, EC2 \texttt{/auth/...} endpoint'leri ile
uçtan uca doğrulanmıştır. Kayıt sonrası otomatik olarak
\texttt{TaxiGuard} BLE çevresel cihazına bağlanmakta, MATCH satırı
geldiğinde \texttt{tel://}<numara> URL'i ile iOS dialer ekranı
otomatik olarak açılmaktadır.

\section{Tartışma}

Sistem mimarisinde alınan üç temel kararın geriye dönüp bakıldığında
hangi pratik faydaları sağladığı şu şekilde değerlendirilebilir:

\begin{itemize}[itemsep=3pt]
    \item \textbf{Kendi sunucumuz vs.\ kapalı kutu API:} İlk bakışta
          AWS Rekognition gibi kapalı kutu bir servis hızlı bir geliştirme
          vaat ediyor olsa da, eşik değerinin kontrolü, model değişiklik
          riski ve KVKK belirsizliği nedeniyle reddedilmiştir. Sergi
          öncesi optimizasyon esnekliği, açık kaynak InsightFace üzerinde
          \texttt{MATCH\_THRESHOLD} ve \texttt{MIN\_QUALITY\_FRAMES}'i
          serbestçe ayarlayabilme imkânını sağlamıştır.
    \item \textbf{Multi-frame vs.\ tek snapshot:} 10 karelik burst, hem
          robust eşleştirme hem ek bir model gerektirmeyen pasif canlılık
          sinyali kazandırmıştır. Maliyeti yalnızca $\sim$5~s'lik
          gecikmedir; bu süre, yolcu binişi başına bir kez gerçekleşen
          bir işlem için makuldür.
    \item \textbf{BLE + telefon vs.\ GSM modülü:} GSM modülü tabanlı
          (SIM800/SIM7600) çözüm, hem donanım maliyeti (en az 250 TL
          ek), hem SIM kart yönetimi hem de operatör servis bedeli gibi
          ekstra yükler getirirdi. Telefonun BLE ile köprüleştirilmesi,
          aynı işlevi sıfır ek altyapıyla sağlamıştır.
    \item \textbf{Entegre BLE vs.\ ayrı modül (HM-10):} İlk Plan B'de yer
          alan HM-10 modülü, ESP32-CAM'in zaten desteklediği BLE
          çevresel arayüzü tarafından yerinden edildi. Bu karar sergi
          öncesinde alındı; tek board üzerinde kamera + Wi-Fi + BLE
          birleşince hem maliyet (~80 TL) hem de ek bir UART hattı
          (STM32 USART2) düştü.
    \item \textbf{iPhone \texttt{tel://} vs.\ Android otomatik çağrı:}
          iOS sandbox tam-otomatik aramaya izin vermez; sürücünün
          dialer ekranında yeşil tuşa tek dokunması istenir
          \cite{apple_tel_url}. Bu, sergi içinde elde olan iPhone~16'nın
          kullanılmasını sağladı ve aynı zamanda yanlış pozitif arama
          koruması olarak savunulabilir bir tasarım sundu.
\end{itemize}

Kısıtlar ve limitler de açıkça raporlanmalıdır:

\begin{itemize}[itemsep=2pt]
    \item Ölçüm seti 5--8 kişilik bir gruba indirgendiği için FAR/FRR
          değerleri istatistiksel olarak sınırlı kapsamlıdır. Sunulan
          metrikler, sistemin \emph{ideal koşullarda yapılabilir} olduğunu
          göstermektedir; geniş ölçekli (n=100+) konuşlandırma için ek
          ölçüm gerekecektir.
    \item Mevcut sunucu HTTP üzerinden çalışmaktadır (Faz 1). Üretim
          ortamı için TLS + token tabanlı kimlik doğrulama zorunludur;
          tasarım buna hazırdır ancak bu tezin değerlendirme kapsamı
          dışındadır.
    \item Yolcunun başörtüsü, gözlük, maske gibi yüz kapayıcı aksesuarlar
          tanımayı etkileyebilmektedir. Bu durum sistem hatası değil,
          ArcFace tabanlı yüz tanımanın doğal bir kısıtıdır ve operasyonel
          süreçte (örnek: birden çok deneme + canlı sürücü onayı) çözülür.
\end{itemize}
